\section{Open questions}

The following five questions were genuinely open at the close of this
replication --- not rhetorical, not answered by the reproduced numerics, and
each with a concrete experimental probe.

\begin{enumerate}
  \item \textbf{MC-VQE vs SA-CASSCF / CASCI at matched cost.} Does MC-VQE
        actually beat classical state-averaged CASSCF or CASCI in a fair
        cost-vs-accuracy comparison on the same exciton systems, or is it only
        competitive with the (weak) CIS baseline the paper contrasts it
        against? \emph{Next step:} reimplement SA-CASSCF and CASCI over the
        same exciton active space (PySCF / OpenFermion), score energy error vs
        FCI and wall-time / operation count vs MC-VQE at matched parameter
        budget across N=8, 12, 18, 24; report accuracy-vs-cost Pareto.

  \item \textbf{MC-VQE vs modern excited-state VQE variants.} How does
        MC-VQE compare quantitatively to SSVQE, VQD, QSE, and MCSCF-VQE on the
        paper's LH2 / H-aggregate benchmarks, and does any one dominate on the
        Pareto? \emph{Next step:} port all four into a common statevector
        framework (Qiskit or PennyLane), run them on the N=8 stack and N=12
        ring at matched entangler depth, and report energy error, oscillator
        error, iteration count, and measurement budget per method.

  \item \textbf{Extension to non-adiabatic dynamics.} Does MC-VQE extend
        cleanly to surface-hopping / Ehrenfest dynamics, which require both
        excited-state energies \emph{and} non-adiabatic coupling vectors
        (NACs) between states? \emph{Next step:} derive an analytic NAC
        estimator from the MC-VQE contracted-Hamiltonian eigenvectors
        (Hammes-Schiffer / Tully overlap formalism), implement, and benchmark
        vs FCI-NACs on a 2-state avoided-crossing (LiF, or a two-chromophore
        exciton at tunable coupling).

  \item \textbf{Noise robustness of state-averaged optimization.} How robust
        is state-averaged MC-VQE to realistic shot noise, readout error, and
        depolarizing noise on today's superconducting hardware, and does the
        state-averaging help or hurt vs single-state VQE under noise?
        \emph{Next step:} add a Qiskit Aer noise model calibrated to a recent
        ibm\_torino or ibm\_kyiv snapshot, run MC-VQE at N=6/8 with $10^3$ to
        $10^6$ shots, and compare vs VQD on the same Hamiltonians. Report
        error scaling in shot count and noise strength.

  \item \textbf{ML-selected active spaces.} Can ML-selected active spaces
        (auto-CAS, orbital-selection neural nets) improve MC-VQE by choosing
        better contracted reference states than CIS, especially for systems
        where CIS is qualitatively wrong (H-aggregates, doubly-excited
        states)? \emph{Next step:} replace CIS reference selection with an
        auto-CAS pipeline (autoCAS) or a small orbital-importance NN, rerun
        on the N=8 stack and on a small doubly-excited test system
        (butadiene $2A_g$), and quantify whether auto-selected references
        reduce required entangler depth or eliminate the doubly-excited-state
        blind spot.
\end{enumerate}
